%%% Fiktivní kapitola s ukázkami tabulek, obrázků a kódu

\chapter{Experimentální část}

V této kapitole popíšeme výsledky provedených experimentů. Nejprve si obecně
popíšeme postup testování, způsob generování vstupních a výstupních souborů a
sledované hodnoty. Následně se budeme věnovat jednotlivým skupinám grafů a
jejich vyhodnocení.

Vstupní data byla vytvořena pomocí skriptu \texttt{geninput.py}, jehož fungování
bylo popsáno v sekci 2.1. Celkem bylo vygenerováno 26 vstupních instancí
odpovídajících problémům popsaným v sekci 1.3.

Generování výstupních SAT formulí probíhalo pomocí programu \texttt{sym.py}.
Pro experimenty byla použita upravená verze tohoto programu, která umožňovala
jednodušší hromadné spouštění programu na stejný vstup s různými parametry. Pro každý
vstup byla použita určitá podmnožina množin permutací popsaných v sekci 2.3 a pro každou zvolenou množinu
permutací bylo následně testováno několik hodnot $k_{\max}$.

Hodnoty $k_{\max}$ byly určeny podle celkového počtu proměnných
v daném zadání. Pokud tento počet označíme $K$, nejprve jsme jej snížili na
největší číslo $K'$ takové, že $K' \leq K$ a $K'$ je dělitelné pěti. Poté byly použity
hodnoty
\[
    \frac{1}{5}K', \frac{2}{5}K', \frac{3}{5}K', \frac{4}{5}K', K'.
\]
Kvůli jednotnému názvu výstupních souborů byla navíc použita také hodnota
$k_{\max}=0$. Program v tomto případě v podstatě 
jen vstup uložil do souboru s odlišným názvem pro jednodušší zpracování.
Celkem tedy bylo pro každou zvolenou dvojici vstupu a množiny permutací uvažováno 
šest hodnot parametru $k_{\max}$.

Protože hodnota $K'$ je ve většině případů menší než skutečný počet proměnných
$K$, nejvyšší testovaná hodnota $k_{\max}$ neodpovídá úplnému průchodu
celého vektoru proměnných. Nejsou tedy nutně vygenerovány všechny klauzule,
které by bylo možné pro danou množinu permutací získat. Z předchozího testování
však vyplynulo, že pro dostatečně velké instance bývají klauzule odpovídající
posledním několika pozicím vektoru téměř vždy již pokryty dříve vygenerovanými
omezeními a další symetrie zpravidla nerozbíjejí. Toto zjednodušení je tedy pro
uvažované instance prakticky přijatelné. U výrazně menších instancí by však takové
omezení mohlo mít znatelnější vliv; například u formule s devíti proměnnými by
ukončení generování již na hodnotě $k_{\max}=5$ mohlo vynechat podstatnou část
potenciálně užitečných klauzulí.

Kromě množin permutací implementovaných přímo v programu \texttt{sym.py} byla
pro všechny problémy testována také množina označená jako \texttt{Satsuma}. Tato
množina obsahuje generátory získané pomocí nástroje \texttt{satsuma}~\cite{DBLP:journals/corr/abs-2406-13557}.
Pro její získání byla použita mírně upravená verze nástroje \texttt{satsuma}, která
v rámci běhu uložila nalezenou množinu permutací do souboru. Takto získaná
množina byla následně použita jako množina permutací pro \texttt{sym.py} namísto některé z množin
permutací generovaných přímo naším nástrojem. Díky tomu můžeme porovnat,
zda množina generátorů nalezená nástrojem \texttt{satsuma} vede k lepšímu rozbití
symetrií než ostatní testované množiny generátorů. Samostatně byly navíc testovány
také výstupní SAT formule vytvořené přímo nástrojem \texttt{satsuma}, který byl spuštěn ve variantě \texttt{satsuma lex ...}. Tento přístup je totiž bližší přístupu našeho nástroje. Druhou možností je \texttt{satsuma fix ...}, která pracuje jinak a zdá se, že dosahuje obecně lepších výsledků. Výsledky první varianty nástroje \texttt{satsuma} budou porovnávány s výsledky našeho nástroje průběžně v dalších sekcích.

Všechny výsledné SAT formule byly testovány pomocí SAT řešiče \texttt{kissat}~\cite{kissat}
a dále pomocí nástroje \texttt{d4}~\cite{d4}. Nástroj \texttt{d4}
umožňuje u splnitelných formulí spočítat celkový počet modelů. Tato hodnota
je pro naše účely velmi důležitá, protože přímo ukazuje, kolik řešení zůstalo po
přidání klauzulí pro rozbíjení symetrií. Na rozdíl od času běhu SAT řešiče, který
může být ovlivněn mnoha dalšími faktory, je počet modelů mnohem lepším
ukazatelem síly rozbití symetrií.

Počítání všech modelů je však výrazně náročnější než samotné rozhodnutí
splnitelnosti. Z tohoto důvodu nebylo možné použít \texttt{d4} pro všechny testované
instance. V mnoha případech výpočet nedoběhl v časovém limitu, který byl nastaven
na 10 hodin. U těchto instancí proto sledujeme především čas běhu řešiče
\texttt{kissat}. Pokud bylo možné použít \texttt{d4}, sledujeme také počet modelů
výsledné formule.

U každého testu byly sledovány především tři hodnoty. První hodnotou je
čas běhu našeho nástroje, tedy čas potřebný k vygenerování klauzulí pro rozbíjení
symetrií. Tento čas ukládal program \texttt{sym.py} přímo do komentářů výstupní
SAT formule. Druhou hodnotou je čas běhu externího nástroje, tedy buď
\texttt{kissat}, nebo \texttt{d4}. Při porovnávání časů je důležité čas potřebný pro samotné vygenerování přidaných klauzulí brát v úvahu, protože pokud
by zrychlení řešiče bylo menší než čas potřebný pro běh našeho nástroje, nemuselo
by být použití nástroje pro danou instanci prakticky výhodné. Poslední hodnotou je počet modelů, pokud jej bylo možné
spočítat.

Budeme porovnávat rozdíl mezi řešením původní
formule a formule po přidání klauzulí pro rozbíjení symetrií, vliv různých hodnot parametru $k_{\max}$, vliv různých množin permutací a zejména
rozdíly mezi generujícími a negenerujícími množinami. Zvláštní pozornost
budeme věnovat také množině permutací získané z nástroje \texttt{satsuma} a
porovnání výsledků našeho nástroje s výsledky nástroje \texttt{satsuma}.

Výsledky jsou v dalších částech uváděny v grafech vygenerovaných pomocným skriptem,
který využíval knihovny \texttt{matplotlib} \cite{Hunter:2007}, \texttt{pandas} \cite{reback2020pandas} a \texttt{numpy} \cite{harris2020array}. 
Zpravidla jsou dvě varianty použitých grafů pro každý problém. Buď je použit graf s časy běhu řešiče \texttt{kissat} společně s grafem s časy běhu našeho nástroje, nebo je použit graf se součty časů běhu \texttt{d4} a našeho nástroje společně s grafem s počty různých modelů problému. V příloze je však pro každé zadání dostupné větší množství grafů obsahujících kromě časů a počtu řešení i poměr zrychlení oproti vstupu, poměr času 
běhu nástroje oproti \texttt{kissat} nebo \texttt{d4}, řadu heatmap a také souhrnné tabulky všech výsledků.

\section{Problém holubníku}

Výsledky pro problém holubníku jsou zobrazeny v grafech na obrázcích \ref{obr11}, \ref{obr12}, \ref{obr13}, \ref{obr14} a \ref{obr15}. Všechny testované instance tohoto problému jsou nesplnitelné. Přestože bychom u splnitelných instancí mohli sledovat počet modelů, standardně se u tohoto problému uvažují jen nesplnitelné instance a navíc splnitelných instancí uvidíme dostatek u ostatních problémů.

Ze všech námi vytvořených testovaných množin permutací dosahuje jednoznačně nejlepších výsledků množina \texttt{Neighbors}. U všech testovaných instancí zde dochází k výraznému zkrácení času běhu SAT řešiče, přičemž čas potřebný pro samotné vygenerování klauzulí se pohybuje nejvýše okolo jedné sekundy. Získané zrychlení je tedy řádově vyšší než režie spojená s během našeho nástroje, takže jeho použití je v tomto případě jednoznačně přínosné.

Naopak většina ostatních testovaných množin permutací nedokáže běh SAT řešiče zrychlit natolik, aby byly náročnější instance vyřešeny před uplynutím nastaveného časového limitu. Pravděpodobným důvodem je, že nesplnitelné instance problému holubníku se stávají pro SAT řešič velmi obtížnými již při relativně malém počtu symetrických holubů a holubníků. Pokud po přidání klauzulí pro rozbíjení symetrií zůstane dostatečně velká symetrická podmnožina řádků nebo sloupců, zůstává problém výpočetně velmi náročný.

Množina \texttt{Neighbors} se od ostatních množin liší tím, že obsahuje transpozice odpovídající záměnám všech sousedních řádků i všech sousedních sloupců matice. Rozbíjení symetrií je tak rovnoměrně rozprostřeno přes celý vektor proměnných a nezůstávají větší souvislé oblasti se zachovanou symetrií. Naproti tomu množiny \texttt{Randtrans} a \texttt{Nongen} sice obsahují stejný počet permutací jako \texttt{Neighbors}, jejich permutace jsou však vybírány náhodně, takže není zaručeno jejich rovnoměrné rozmístění mezi všemi proměnnými. Přesto právě tyto dvě množiny dosahují po množině \texttt{Neighbors} nejlepších výsledků, což naznačuje, že samotný počet použitých permutací hraje významnou roli.

Množina \texttt{Transpositions} se pro uvažované velikosti instancí ukázala jako prakticky nepoužitelná, protože pro rozměry binárních matic uvažované v uvedených problémech byl nárůst času potřebného ke generování klauzulí tak vysoký, že buď ve většině případů běh nástroje ani neskončil v rozumném čase, nebo by oproti množině \texttt{Neighbors} byl velmi nevýhodný. Pro problém holubníku uvažujeme množinu \texttt{Transpositions} zvlášť pro řádky a pro sloupce a poté procházíme jejich kartézský součin postupem popsaným v sekci 1.2. To znamená, že pro $pig$ řádků máme celkem $\frac{pig(pig-1)}{2}$ permutací a pro $hol$ sloupců celkem $\frac{hol(hol-1)}{2}$ permutací. Ke každé množině přidáme identitu a poté procházíme všechny možné dvojice prvků z těchto množin (kromě dvojice, ve které je na obou pozicích příslušná identická permutace). Celkem dostáváme $(\frac{pig(pig-1)}{2}+1)(\frac{hol(hol-1)}{2}+1)-1$ různých permutovaných vektorů proměnných. Například pro \texttt{php\_pig31\_hol30} dostáváme tedy $(\frac{31(31-1)}{2}+1)(\frac{30(30-1)}{2}+1)-1=203175$. Množina \texttt{Neighbors} v tomto případě pro řádky obsahuje $30$ permutací a pro sloupce $29$ permutací. Analogickým postupem se dostaneme k celkovému počtu $929$ různých permutovaných vektorů proměnných. Jak je vidět, rozdíl mezi počty generovaných permutovaných vektorů proměnných je více než dvousetnásobný. Není proto překvapivé, že se tato skutečnost výrazně projeví na době běhu našeho nástroje. Tento rozdíl navíc s rostoucí velikostí problému dále narůstá, protože počet permutací v množině \texttt{Transpositions} roste kvadraticky v počtu řádků i sloupců ($\Theta(pig^2\cdot hol^2)$), zatímco u množiny \texttt{Neighbors} pouze lineárně ($\Theta(pig\cdot hol)$).

Na instanci \texttt{php\_pig40\_hol10} lze pozorovat zajímavý jev, který se bude opakovat i u dalších problémů. Nejkratší čas běhu SAT řešiče zde není dosažen pro nejvyšší hodnotu $k_{\max}$, ale pro některou z nižších hodnot. Přestože vyšší hodnoty $k_{\max}$ vždy rozbijí alespoň tolik symetrií jako nižší hodnoty, zároveň vedou k vytvoření většího množství pomocných proměnných a klauzulí. V některých případech může být režie spojená s jejich zpracováním větší než přínos více rozbitých symetrií. Při hledání vhodné hodnoty parametru $k_{\max}$ je proto nutné zohlednit nejen účinnost rozbití symetrií a čas generování klauzulí, ale také vliv vyššího množství proměnných na výkon SAT řešiče.

Kromě \texttt{Neighbors} je jedinou výrazně účinnější množinou permutací množina \texttt{Satsuma}, která \texttt{Neighbors} dokonce ve většině případů překonává. Co je na tom nejvíce překvapivé je to, že množina \texttt{Satsuma} obsahuje výrazně méně permutací než množina \texttt{Neighbors}. Při předchozím testování se totiž ukázalo, že celkový počet použitých permutací má obecně významný vliv na účinnost rozbití symetrií. Zde se celkový počet permutací v množině \texttt{Satsuma} pro testované instance pohybuje v rozmezí od $40$ do $59$, zatímco v případě \texttt{Neighbors} se používá kartézský součin transpozic pro řádky a pro sloupce. To znamená celkem $pig\cdot hol-1$ permutací, což se zde pohybuje v řádu několika stovek pro všechny instance. Permutace nalezené při detekci symetrií nástrojem \texttt{satsuma} tedy v tomto případě opravdu dávají velmi dobré částečné rozbití symetrie, jak bylo zmíněno v první kapitole. Jedním z možných vysvětlení je skutečnost, že problém holubníku patří mezi nejčastěji používané benchmarky při výzkumu rozbíjení symetrií, a i nástroje pro detekci symetrií proto mohou být na podobné struktury obzvláště dobře přizpůsobeny. 

Nakonec pokud srovnáme nástroj \texttt{satsuma} a náš nástroj s množinou \texttt{Satsuma}, dostaneme téměř totožné výsledky. SAT řešič v obou případech zjistí téměř okamžitě, že je problém nesplnitelný.

\begin{figure}[H]
\centering
\begin{subfigure}{0.48\textwidth}
        \centering
        \includegraphics[width=\textwidth]{img/plots/kissat_time_log/php_pig21_hol20.png}
    \end{subfigure}
    \hfill
    \begin{subfigure}{0.48\textwidth}
        \centering
        \includegraphics[width=\textwidth]{img/plots/tool_time_log/php_pig21_hol20.png}
    \end{subfigure}
\caption{Čas kissat a čas nástroje pro php\_pig21\_hol20}
\label{obr11}
\end{figure}

\begin{figure}[H]
\centering
\begin{subfigure}{0.48\textwidth}
        \centering
        \includegraphics[width=\textwidth]{img/plots/kissat_time_log/php_pig25_hol20.png}
    \end{subfigure}
    \hfill
    \begin{subfigure}{0.48\textwidth}
        \centering
        \includegraphics[width=\textwidth]{img/plots/tool_time_log/php_pig25_hol20.png}
    \end{subfigure}
\caption{Čas kissat a čas nástroje pro php\_pig25\_hol20}
\label{obr12}
\end{figure}

\begin{figure}[H]
\centering
\begin{subfigure}{0.48\textwidth}
        \centering
        \includegraphics[width=\textwidth]{img/plots/kissat_time_log/php_pig31_hol30.png}
    \end{subfigure}
    \hfill
    \begin{subfigure}{0.48\textwidth}
        \centering
        \includegraphics[width=\textwidth]{img/plots/tool_time_log/php_pig31_hol30.png}
    \end{subfigure}
\caption{Čas kissat a čas nástroje pro php\_pig31\_hol30}
\label{obr13}
\end{figure}

\begin{figure}[H]
\centering
\begin{subfigure}{0.48\textwidth}
        \centering
        \includegraphics[width=\textwidth]{img/plots/kissat_time_log/php_pig35_hol15.png}
    \end{subfigure}
    \hfill
    \begin{subfigure}{0.48\textwidth}
        \centering
        \includegraphics[width=\textwidth]{img/plots/tool_time_log/php_pig35_hol15.png}
    \end{subfigure}
\caption{Čas kissat a čas nástroje pro php\_pig35\_hol15}
\label{obr14}
\end{figure}

\begin{figure}[H]
\centering
\begin{subfigure}{0.48\textwidth}
        \centering
        \includegraphics[width=\textwidth]{img/plots/kissat_time_log/php_pig40_hol10.png}
    \end{subfigure}
    \hfill
    \begin{subfigure}{0.48\textwidth}
        \centering
        \includegraphics[width=\textwidth]{img/plots/tool_time_log/php_pig40_hol10.png}
    \end{subfigure}
\caption{Čas kissat a čas nástroje pro php\_pig40\_hol10}
\label{obr15}
\end{figure}


\section{Problém N královen}

Výsledky pro problém N královen jsou zobrazeny na obrázcích \ref{obr16}, \ref{obr17}, \ref{obr18}, \ref{obr19} a \ref{obr20}. Na rozdíl od problému holubníku jsou všechny testované instance splnitelné. Přestože nalezení jednoho řešení je díky vysokému počtu různých řešení pro SAT řešič poměrně jednoduché, je zajímavé sledovat právě celkový počet modelů nalezený nástrojem \texttt{d4}.

Pro tento problém porovnáváme pouze naši generující množinu permutací dihedrální grupy \texttt{DihedralGen} a množinu \texttt{Satsuma}. Množina \texttt{DihedralGen} obsahuje vždy dvě permutace, zatímco množina \texttt{Satsuma} obsahuje dvě nebo tři permutace. Množina obsahuje tři permutace v případech $n=10$ a $n=13$. Právě v těchto případech můžeme v grafech pozorovat nejvýraznější rozdíly mezi testovanými množinami permutací. U ostatních instancí, kde jsou množiny stejně velké, již rozdíly nejsou tak výrazné. Toto pozorování potvrzuje významný vliv počtu použitých permutací na účinnost rozbíjení symetrií, který byl zmíněn v předchozí sekci.

Problém N královen se od ostatních problémů liší také tím, že počet jeho symetrií nezávisí na velikosti instance. Pro všechna $n$ jsou symetrie dány osmi symetriemi čtvercové šachovnice tvořícími dihedrální grupu. Se zvyšující se velikostí šachovnice se tedy zvětšuje pouze SAT formule, zatímco množina symetrií zůstává stále stejná.

V porovnání s nástrojem \texttt{satsuma} jsou výsledky našeho nástroje podobné. Výraznější rozdíly lze pozorovat pouze u instancí, u kterých nástroj \texttt{satsuma} při vlastním běhu nalezl pravděpodobně množinu o třech permutacích, zatímco při exportu množiny byly nalezeny pouze dvě permutace. Tyto dva běhy totiž nebyly prováděny současně a nástroj může nalézt pro totožné zadání různé množiny generátorů pro různé běhy. 

Časy běhu nástroje \texttt{d4} jsou u většiny testovaných instancí velmi malé, a proto z nich nelze vyvozovat jednoznačné závěry. Obecně nejvíce vypovídajícím se jeví případ s nejvyšším $n$, tedy $n=14$. Zde už se časy běhu pohybují v řádu minut a lze jednoznačně vidět, že pro množinu \texttt{DihedralGen} se čas významně sníží pro první nenulovou hodnotu $k_{max}$ a poté postupně roste. Toto chování odpovídá také grafům s počty modelů. Dramatický pokles totiž nastane právě pro první nenulovou hodnotu $k_{max}$ a poté už klesá jen velmi mírně. Tím pádem pro vyšší $k_{max}$ už téměř žádné zrychlení rozbíjení symetrií nepřináší a přidávání dalších pomocných proměnných jen řešení zpomaluje. Proto čas roste.

Zajímavou vlastností grafů s počtem modelů je jejich konzistence napříč různými velikostmi problému. Pro množinu \texttt{DihedralGen}, která obsahuje vždy stejné dva generátory dihedrální grupy, se poměr počtu modelů po rozbití symetrií vůči původnímu počtu modelů ve všech testovaných instancích pohybuje lehce pod hodnotou $\frac{4}{10}$. Stejně konzistentně se projevuje i závislost na parametru $k_{max}$, kdy největší pokles počtu modelů nastává již pro první nenulovou hodnotu tohoto parametru. U množiny \texttt{Satsuma} lze pozorovat, že výraznější poklesy počtu modelů nastávají v různých instancích při různých hodnotách $k_{max}$, což naznačuje, že optimální hloubka generování by mohla záviset na konkrétních použitých permutacích. U množiny \texttt{Satsuma} se totiž narozdíl od množiny \texttt{DihedralGen} mohou pro různé běhy vyskytovat různé generátory dihedrální grupy. Při dalším testování by tedy bylo žádoucí zahrnout všech osm permutací a sledovat, pro které hodnoty $k_{max}$ se jednotlivé permutace projeví nejvýrazněji a jestli jsou hodnoty podobně konzistentní napříč různými velikostmi zadání.

\begin{figure}[H]
\centering
\begin{subfigure}{0.48\textwidth}
        \centering
        \includegraphics[width=\textwidth]{img/plots/total_d4_time_log/queens_n10.png}
    \end{subfigure}
    \hfill
    \begin{subfigure}{0.48\textwidth}
        \centering
        \includegraphics[width=\textwidth]{img/plots/models_ratio_log/queens_n10.png}
    \end{subfigure}
\caption{Čas nástroje + čas d4 a počet modelů pro queens\_n10}
\label{obr16}
\end{figure}

\begin{figure}[H]
\centering
\begin{subfigure}{0.48\textwidth}
        \centering
        \includegraphics[width=\textwidth]{img/plots/total_d4_time_log/queens_n11.png}
    \end{subfigure}
    \hfill
    \begin{subfigure}{0.48\textwidth}
        \centering
        \includegraphics[width=\textwidth]{img/plots/models_ratio_log/queens_n11.png}
    \end{subfigure}
\caption{Čas nástroje + čas d4 a počet modelů pro queens\_n11}
\label{obr17}
\end{figure}

\begin{figure}[H]
\centering
\begin{subfigure}{0.48\textwidth}
        \centering
        \includegraphics[width=\textwidth]{img/plots/total_d4_time_log/queens_n12.png}
    \end{subfigure}
    \hfill
    \begin{subfigure}{0.48\textwidth}
        \centering
        \includegraphics[width=\textwidth]{img/plots/models_ratio_log/queens_n12.png}
    \end{subfigure}
\caption{Čas nástroje + čas d4 a počet modelů pro queens\_n12}
\label{obr18}
\end{figure}

\begin{figure}[H]
\centering
\begin{subfigure}{0.48\textwidth}
        \centering
        \includegraphics[width=\textwidth]{img/plots/total_d4_time_log/queens_n13.png}
    \end{subfigure}
    \hfill
    \begin{subfigure}{0.48\textwidth}
        \centering
        \includegraphics[width=\textwidth]{img/plots/models_ratio_log/queens_n13.png}
    \end{subfigure}
\caption{Čas nástroje + čas d4 a počet modelů pro queens\_n13}
\label{obr19}
\end{figure}

\begin{figure}[H]
\centering
\begin{subfigure}{0.48\textwidth}
        \centering
        \includegraphics[width=\textwidth]{img/plots/total_d4_time_log/queens_n14.png}
    \end{subfigure}
    \hfill
    \begin{subfigure}{0.48\textwidth}
        \centering
        \includegraphics[width=\textwidth]{img/plots/models_ratio_log/queens_n14.png}
    \end{subfigure}
\caption{Čas nástroje + čas d4 a počet modelů pro queens\_n14}
\label{obr20}
\end{figure}

\section{Problém nalezení matice s řádky a sloupci se zadanými součty}

Výsledky experimentů pro problém nalezení matice s řádky a sloupci se zadanými součty jsou uvedeny na obrázcích \ref{obr21}, \ref{obr23}, \ref{obr25} a \ref{obr26}. Testovány byly také instance s parametry $rc=10, rcrand=5$ a $rc=13, rcrand=10$. Jelikož však jejich výsledky byly velmi podobné instanci s parametry $rc=15, rcrand=3$, nejsou z důvodu úspory místa v hlavním textu uvedeny. Jejich grafy jsou dostupné v příloze práce.

Porovnávat v tomto případě můžeme opět pouze dvě množiny permutací a to množinu \texttt{SumsNeighbors} a množinu \texttt{Satsuma}. U instancí na obrázcích \ref{obr21} a \ref{obr23} lze pozorovat velmi podobné chování jako u množiny \texttt{Neighbors} v problému holubníku. Dochází zde k výraznému snížení počtu modelů i času běhu řešiče a účinnost rozbíjení symetrií se stále zvyšuje i pro vyšší hodnoty parametru $k_{max}$. Podobnost není náhodná. Pokud uvažujeme matici o $11$ řádcích a $10$ sloupcích, ve které mají všechny řádky i všechny sloupce předepsaný součet roven jedné, dostáváme téměř stejný problém jako je problém holubníku s $11$ holuby a $10$ holubníky. 

U ostatních testovaných instancí již experimenty neposkytují tak jednoznačné závěry. Ve splnitelných instancích je počet modelů natolik vysoký, že nástroj \texttt{d4} nedokončí jejich spočítání v daném časovém limitu, zatímco SAT řešič \texttt{kissat} nalezne jedno řešení téměř okamžitě. Naměřené časy řešení se proto pohybují v řádu milisekund a jejich vzájemné rozdíly jsou velmi malé. Tyto rozdíly mohou být ovlivněny řadou dalších faktorů uvnitř SAT řešiče a nelze z nich proto vyvozovat významnější závěry o kvalitě jednotlivých množin permutací. Obdobná situace nastává u nesplnitelné instance na obrázku \ref{obr25}, kde je pravděpodobné, že tato instance obsahuje poměrně jednoduchý spor vyplývající ze zvolených hodnot parametrů $rc$ a $rcrand$, takže SAT řešič dokáže nesplnitelnost prokázat téměř okamžitě a opět čas ovlivňují spíše náhodné vlivy běhu SAT řešiče.

Při porovnání množiny \texttt{Satsuma} a samotného nástroje \texttt{satsuma} dostáváme pro tento problém odlišné výsledky než u předchozích experimentů. Množina \texttt{Satsuma} obsahuje podstatně méně permutací než množina \texttt{SumsNeighbors}, což se projevuje nižší účinností rozbití symetrií. Ani samotný nástroj \texttt{satsuma} zde nedosahuje výsledků srovnatelných s naším nástrojem využívajícím množinu \texttt{SumsNeighbors}. U instancí na obrázcích \ref{obr21} a \ref{obr23} jsou výsledky nástroje \texttt{satsuma} pouze mírně lepší než výsledky našeho nástroje používajícího exportovanou množinu \texttt{Satsuma}, zatímco množina \texttt{SumsNeighbors} vede k výrazně lepším výsledkům než oba tyto přístupy.

Možným vysvětlením je charakter symetrií tohoto problému. Množina permutací \texttt{SumsNeighbors} je navržena přímo pro symetrie vznikající mezi řádky nebo sloupci se stejnými součty, zatímco nástroj \texttt{satsuma} detekuje symetrie pouze ze zadané SAT formule. Zdá se tedy, že symetrie pouze na podmnožinách řádků nebo sloupců jsou pro nástroj \texttt{satsuma} těžko detekovatelné.

\begin{figure}[H]
\centering
\begin{subfigure}{0.48\textwidth}
        \centering
        \includegraphics[width=\textwidth]{img/plots/total_d4_time_log/rcsums_rc10_rcrand2.png}
    \end{subfigure}
    \hfill
    \begin{subfigure}{0.48\textwidth}
        \centering
        \includegraphics[width=\textwidth]{img/plots/models_ratio_log/rcsums_rc10_rcrand2.png}
    \end{subfigure}
\caption{Čas nástroje + čas d4 a počet modelů pro rcsums\_rc10\_rcrand2}
\label{obr21}
\end{figure}

%\begin{figure}[H]
%\centering
%\begin{subfigure}{0.48\textwidth}
%        \centering
%        \includegraphics[width=\textwidth]{img/plots/kissat_time_log/rcsums_rc10_rcrand5.png}
%    \end{subfigure}
%    \hfill
%    \begin{subfigure}{0.48\textwidth}
%        \centering
%        \includegraphics[width=\textwidth]{img/plots/tool_time_log/rcsums_rc10_rcrand5.png}
%    \end{subfigure}
%\caption{Celkový čas řešení pomocí}
%\label{obr22:rcsums-rc10-rcrand5-total}
%\end{figure}

\begin{figure}[H]
\centering
\begin{subfigure}{0.48\textwidth}
        \centering
        \includegraphics[width=\textwidth]{img/plots/total_d4_time_log/rcsums_rc12_rcrand2.png}
    \end{subfigure}
    \hfill
    \begin{subfigure}{0.48\textwidth}
        \centering
        \includegraphics[width=\textwidth]{img/plots/models_ratio_log/rcsums_rc12_rcrand2.png}
    \end{subfigure}
\caption{Čas nástroje + čas d4 a počet modelů pro rcsums\_rc12\_rcrand2}
\label{obr23}
\end{figure}

%\begin{figure}[H]
%\centering
%\begin{subfigure}{0.48\textwidth}
%        \centering
%        \includegraphics[width=\textwidth]{img/plots/kissat_time_log/rcsums_rc13_rcrand10.png}
%    \end{subfigure}
%    \hfill
%    \begin{subfigure}{0.48\textwidth}
%        \centering
%        \includegraphics[width=\textwidth]{img/plots/tool_time_log/rcsums_rc13_rcrand10.png}
%    \end{subfigure}
%\caption{Celkový čas řešení pomocí}
%\label{obr24:rcsums-rc13-rcrand10-total}
%\end{figure}

\begin{figure}[H]
\centering
\begin{subfigure}{0.48\textwidth}
        \centering
        \includegraphics[width=\textwidth]{img/plots/kissat_time_log/rcsums_rc14_rcrand13.png}
    \end{subfigure}
    \hfill
    \begin{subfigure}{0.48\textwidth}
        \centering
        \includegraphics[width=\textwidth]{img/plots/tool_time_log/rcsums_rc14_rcrand13.png}
    \end{subfigure}
\caption{Čas kissat a čas nástroje pro rcsums\_rc14\_rcrand13}
\label{obr25}
\end{figure}

\begin{figure}[H]
\centering
\begin{subfigure}{0.48\textwidth}
        \centering
        \includegraphics[width=\textwidth]{img/plots/kissat_time_log/rcsums_rc15_rcrand3.png}
    \end{subfigure}
    \hfill
    \begin{subfigure}{0.48\textwidth}
        \centering
        \includegraphics[width=\textwidth]{img/plots/tool_time_log/rcsums_rc15_rcrand3.png}
    \end{subfigure}
\caption{Čas kissat a čas nástroje pro rcsums\_rc15\_rcrand3}
\label{obr26}
\end{figure}



\section{Problém latinských čtverců}

Výsledky pro problém latinských čtverců jsou uvedeny na obrázcích \ref{obr06}, \ref{obr07}, \ref{obr08} a \ref{obr09}. Jedná se o první testovaný problém o třech rozměrech, kterými jsou řádky, sloupce a hodnoty symbolů, které do matice rozmisťujeme. Z důvodu úspory místa nejsou uvedeny výsledky pro instanci $lsq=9$, které jsou velmi podobné instanci $lsq=8$. Grafy pro tuto instanci jsou opět uvedeny v příloze práce.

U dvou nejmenších instancí lze porovnat dosud největší počet různých množin permutací. Nejprve se zaměřme na množiny \texttt{Neighbors}, \texttt{Randtrans} a \texttt{Nongen}. Rozdíly u těchto množin jsme mohli zatím sledovat pouze omezeně u problému holubníku a chceme na nich pozorovat, zda skutečnost, že množina permutací tvoří generující množinu grupy symetrií, vede také k efektivnějšímu rozbíjení symetrií. Ve všech testovaných případech dosahuje opět zdaleka nejlepších výsledků množina \texttt{Neighbors}. Ta využívá transpozice sousedních prvků rovnoměrně rozložené přes všechny tři rozměry problému. Naproti tomu množiny \texttt{Randtrans} a \texttt{Nongen} využívají náhodně vybrané transpozice. Stejně jako u problému holubníku se tak potvrzuje výhoda rovnoměrného pokrytí proměnných.

Zajímavější je vzájemné porovnání množin \texttt{Randtrans} a \texttt{Nongen}. Obě množiny obsahují stejný počet permutací a liší se především tím, že \texttt{Randtrans} může tvořit generující množinu příslušné grupy symetrií, zatímco \texttt{Nongen} nikoliv. Ve všech testovaných případech kromě jediné výjimky pro $lsq = 5$ a $k_{max}=24$ dosahuje množina \texttt{Randtrans} lepších výsledků jak z hlediska počtu modelů, tak i času běhu řešiče. Stejný závěr bylo možné pozorovat také v několika předchozích sériích experimentů. Výsledky proto naznačují, že vlastnost být generující množinou může mít pozitivní vliv na účinnost rozbíjení symetrií. Vyplývá z toho, že by mohlo mít smysl v budoucnu tuto vlastnost studovat i teoreticky a pokusit se ji dokázat, nebo případně vyvrátit, pokud na výsledky testů mají vliv jiné nyní neznámé faktory.

Další zajímavé srovnání představují množiny \texttt{Gen2} a \texttt{Satsuma}. Množina permutací \texttt{Satsuma} obsahuje ve všech testovaných instancích pět permutací, zatímco množina \texttt{Gen2} pouze dvě. Samotný počet permutací však v tomto případě nevystihuje skutečný počet použitých permutovaných vektorů proměnných. Permutace množiny \texttt{Satsuma} jsou definovány přímo na celém vektoru proměnných, zatímco permutace množiny \texttt{Gen2} jsou definovány samostatně pro řádky, sloupce a hodnoty symbolů. Náš nástroj proto vytváří všechny jejich možné kombinace. Trojice mohou na všech pozicích kromě permutací z \texttt{Gen2} také obsahovat identickou permutaci. Celkem tak využívá $3^3-1=26$ různých permutací, přičemž jedinou vynechanou kombinací je trojice identických permutací na všech třech pozicích. Skutečný počet použitých permutací je tedy i v tomto případě vyšší než u množiny \texttt{Satsuma}, což se odráží také na výsledcích.

Nástroj \texttt{satsuma} i v tomto případě dosahuje velmi podobných výsledků jako náš nástroj využívající množinu \texttt{Satsuma}. To naznačuje, že navržený způsob generování Lex-Leader klauzulí dokáže nalezenou množinu generátorů využít srovnatelně efektivně jako nástroj \texttt{satsuma}. Výsledky v grafech pro množinu \texttt{Satsuma} tedy můžeme zároveň považovat za výsledky nástroje \texttt{satsuma}. Ve srovnání s množinou \texttt{Neighbors} však náš nástroj dosahuje lepších výsledků z hlediska redukce počtu modelů. 

Za zmínku stojí také množina \texttt{Transpositions}. Přestože její použití vede kvůli velkému počtu permutací k výrazně delším časům generování i řešení, z hlediska samotného počtu odstraněných symetrií dosahuje ještě lepších výsledků než množina \texttt{Neighbors}. Počet modelů se zde snižuje až na hodnoty $5$, respektive $252$, pro instance $lsq=5$, respektive $lsq=6$. Zároveň je to jeden z mála případů, kdy se na celkovém čase významněji podílí čas běhu nástroje. Proto množinu \texttt{Transpositions} nelze použít pro $lsq>6$. Již pro $lsq=7$ nebylo možné vygenerovat klauzule pro všechny používané hodnoty $k_{max}$. Navíc se kromě času dramaticky zvyšovala i velikost výsledného souboru. To bylo mimochodem limitujícím faktorem i pro řadu dalších instancí u všech uvažovaných problémů.

Pro $lsq=8$ (a $lsq=9$) nastane již známá situace, kdy počet řešení je tak veliký, že běh nástroje \texttt{d4} ani po přidání klauzulí pro rozbíjení symetrií neskončí v časovém limitu, ale naopak nástroj \texttt{kissat} díky velkému počtu řešení nalezne jedno řešení téměř okamžitě. Proto časy kolísají a spíše se zhoršují s přidáváním většího množství pomocných proměnných. Díky nižšímu počtu permutací zde pro množinu \texttt{Satsuma} čas stoupá pomaleji, protože není přidáváno takové množství pomocných proměnných.

\begin{figure}[H]
\centering
\begin{subfigure}{0.48\textwidth}
        \centering
        \includegraphics[width=\textwidth]{img/plots/total_d4_time_log/latinsq_lsq5.png}
    \end{subfigure}
    \hfill
    \begin{subfigure}{0.48\textwidth}
        \centering
        \includegraphics[width=\textwidth]{img/plots/models_ratio_log/latinsq_lsq5.png}
    \end{subfigure}
\caption{Čas nástroje + čas d4 a počet modelů pro golf\_g3\_p3\_w3}
\label{obr06}
\end{figure}

\begin{figure}[H]
\centering
\begin{subfigure}{0.48\textwidth}
        \centering
        \includegraphics[width=\textwidth]{img/plots/total_d4_time_log/latinsq_lsq6.png}
    \end{subfigure}
    \hfill
    \begin{subfigure}{0.48\textwidth}
        \centering
        \includegraphics[width=\textwidth]{img/plots/models_log/latinsq_lsq6.png}
    \end{subfigure}
\caption{Čas nástroje + čas d4 a počet modelů pro latinsq\_lsq6}
\label{obr07}
\end{figure}

\begin{figure}[H]
\centering
\begin{subfigure}{0.48\textwidth}
        \centering
        \includegraphics[width=\textwidth]{img/plots/total_d4_time_log/latinsq_lsq7.png}
    \end{subfigure}
    \hfill
    \begin{subfigure}{0.48\textwidth}
        \centering
        \includegraphics[width=\textwidth]{img/plots/models_log/latinsq_lsq7.png}
    \end{subfigure}
\caption{Čas nástroje + čas d4 a počet modelů pro latinsq\_lsq7}
\label{obr08}
\end{figure}

\begin{figure}[H]
\centering
\begin{subfigure}{0.48\textwidth}
        \centering
        \includegraphics[width=\textwidth]{img/plots/kissat_time_log/latinsq_lsq8.png}
    \end{subfigure}
    \hfill
    \begin{subfigure}{0.48\textwidth}
        \centering
        \includegraphics[width=\textwidth]{img/plots/tool_time_log/latinsq_lsq8.png}
    \end{subfigure}
\caption{Čas kissat a čas nástroje pro latinsq\_lsq8}
\label{obr09}
\end{figure}

%\begin{figure}[H]
%\centering
%\begin{subfigure}{0.48\textwidth}
%        \centering
%        \includegraphics[width=\textwidth]{img/plots/kissat_time_log/latinsq_lsq9.png}
%    \end{subfigure}
%    \hfill
%    \begin{subfigure}{0.48\textwidth}
%        \centering
%        \includegraphics[width=\textwidth]{img/plots/tool_time_log/latinsq_lsq9.png}
%    \end{subfigure}
%\caption{Celkový čas řešení pomocí}
%\label{obr10:latinsq-lsq9-total}
%\end{figure}



\section{Problém společenských golfistů}

Výsledky pro problém společenských golfistů jsou uvedeny na obrázcích \ref{obr01}, \ref{obr02} a \ref{obr03}. Pro přehlednost nejsou zobrazeny grafy pro instance \texttt{golf\_g5\_p4\_w4} a \texttt{golf\_g6\_p2\_w3}, které jsou stejně jako u předchozích problémů uvedeny v příloze práce. Jejich výsledky jsou velmi podobné instanci \texttt{golf\_g5\_p3\_w6}. Tyto instance mají podobně jako už několikrát u předchozích problémů příliš mnoho řešení a tím pádem nástroj \texttt{d4} neskončí v časovém limitu, ale nástroj \texttt{kissat} z velkého množství řešení najde jedno velmi rychle.

Zajímavější jsou proto výsledky na obrázcích \ref{obr01} a \ref{obr02}, kde jedna z testovaných instancí je splnitelná a druhá nesplnitelná. Také zde dosahuje nejlepších výsledků množina \texttt{Neighbors}, což je v souladu s výsledky všech předchozích experimentů. Zajímavé je porovnání množin \texttt{Randtrans} a \texttt{Nongen}. V několika případech pro nižší hodnoty parametru $k_{max}$ dosahuje množina \texttt{Nongen} mírně lepších výsledků než \texttt{Randtrans}. Přestože jsou obě množiny permutací voleny náhodně, takže bychom mohli v krajních případech očekávat občas i výrazné vítězství množiny \texttt{Nongen}, nastává podobná situace velmi ojediněle. Pro vyšší hodnoty $k_{max}$ již množina \texttt{Randtrans} opět dosahuje lepších výsledků, takže i tento experiment podporuje hypotézu, že generující množiny permutací mohou být pro rozbíjení symetrií obecně vhodnější než stejně velké negenerující množiny.

Podobně jako u problému latinských čtverců lze porovnat také množiny \texttt{Gen2} a \texttt{Satsuma}, kdy množina \texttt{Gen2} opět odpovídá celkem $26$ použitým permutacím, zatímco množina \texttt{Satsuma} obsahuje $18$ permutací pro instanci \texttt{golf\_g3\_p3\_w3} a $13$ permutací pro instanci \texttt{golf\_g4\_p4\_w6}. Přesto jsou výsledky obou množin velmi podobné. To naznačuje, že množiny generátorů nalezené nástrojem \texttt{satsuma} při detekci symetrií mají vyšší kvalitu a dokážou i při menším počtu permutací rozbíjet symetrie účinně.

Pro nástroj \texttt{satsuma} je například u \texttt{golf\_g3\_p3\_w3} čas i počet modelů zhruba poloviční oproti našemu nástroji s množinou \texttt{Satsuma}, což ale může být způsobeno rozdílnou nalezenou množinou permutací. Počet modelů po přidání klauzulí je zde $44106$. Pro srovnání náš nástroj s množinou \texttt{Neighbors} se dostal na $12$ modelů. Stejně tak u nesplnitelné instance \texttt{golf\_g4\_p4\_w6} jsou početnější množiny výrazně účinnější.

\begin{figure}[H]
\centering
\begin{subfigure}{0.48\textwidth}
        \centering
        \includegraphics[width=\textwidth]{img/plots/total_d4_time_log/golf_g3_p3_w3.png}
    \end{subfigure}
    \hfill
    \begin{subfigure}{0.48\textwidth}
        \centering
        \includegraphics[width=\textwidth]{img/plots/models_ratio_log/golf_g3_p3_w3.png}
    \end{subfigure}
\caption{Čas nástroje + čas d4 a počet modelů pro golf\_g3\_p3\_w3}
\label{obr01}
\end{figure}

\begin{figure}[H]
\centering
\begin{subfigure}{0.48\textwidth}
        \centering
        \includegraphics[width=\textwidth]{img/plots/kissat_time_log/golf_g4_p4_w6.png}
    \end{subfigure}
    \hfill
    \begin{subfigure}{0.48\textwidth}
        \centering
        \includegraphics[width=\textwidth]{img/plots/tool_time_log/golf_g4_p4_w6.png}
    \end{subfigure}
\caption{Čas kissat a čas nástroje pro golf\_g4\_p4\_w6}
\label{obr02}
\end{figure}

\begin{figure}[H]
\centering
\begin{subfigure}{0.48\textwidth}
        \centering
        \includegraphics[width=\textwidth]{img/plots/kissat_time_log/golf_g5_p3_w6.png}
    \end{subfigure}
    \hfill
    \begin{subfigure}{0.48\textwidth}
        \centering
        \includegraphics[width=\textwidth]{img/plots/tool_time_log/golf_g5_p3_w6.png}
    \end{subfigure}
\caption{Čas kissat a čas nástroje pro golf\_g5\_p3\_w6}
\label{obr03}
\end{figure}

%\begin{figure}[H]
%\centering
%\begin{subfigure}{0.48\textwidth}
%        \centering
%        \includegraphics[width=\textwidth]{img/plots/kissat_time_log/golf_g5_p4_w4.png}
%    \end{subfigure}
%    \hfill
%    \begin{subfigure}{0.48\textwidth}
%        \centering
%        \includegraphics[width=\textwidth]{img/plots/tool_time_log/golf_g5_p4_w4.png}
%    \end{subfigure}
%\caption{Celkový čas řešení pomocí}
%\label{obr04:golf-g5-p4-w4-total}
%\end{figure}

%\begin{figure}[H]
%\centering
%\begin{subfigure}{0.48\textwidth}
%        \centering
%        \includegraphics[width=\textwidth]{img/plots/kissat_time_log/golf_g6_p2_w3.png}
%    \end{subfigure}
%    \hfill
%    \begin{subfigure}{0.48\textwidth}
%        \centering
%        \includegraphics[width=\textwidth]{img/plots/tool_time_log/golf_g6_p2_w3.png}
%    \end{subfigure}
%\caption{Celkový čas řešení pomocí}
%\label{obr05:golf-g6-p2-w3-total}
%\end{figure}

\section{Shrnutí experimentálních výsledků}
\label{sec36}
Provedené experimenty ukázaly, že navržený nástroj je schopen významně urychlit řešení řady SAT formulí pomocí rozbíjení symetrií. Největší přínos lze pozorovat především u obtížných nesplnitelných instancí, kde redukce symetrických řešení vede k výraznému zkrácení času běhu SAT řešiče. Naopak u úloh s velmi vysokým počtem řešení bývá nalezení jednoho modelu i bez rozbíjení symetrií velmi rychlé, a proto se zrychlení běhu SAT řešiče projeví jen omezeně. V těchto případech se jako vhodnější ukázalo sledování celkového počtu modelů, který lépe vystihuje účinnost rozbití symetrií. Přestože počet modelů zde slouží především k hodnocení účinnosti rozbití symetrií, enumerace modelů představuje i samostatnou úlohu. V řadě problémů je totiž cílem nalézt všechna různá řešení daného problému, přičemž započítávat symetrická řešení není žádoucí. V takových případech rozbíjení symetrií neslouží pouze ke zrychlení výpočtu, ale je přímo nezbytnou součástí řešení daného problému.

Ve všech testovaných problémech se jako nejefektivnější z implementovaných množin permutací ukázala množina \texttt{Neighbors}, pokud samozřejmě je na daný problém aplikovatelná. Tato množina rovnoměrně pokrývá celý prostor proměnných pomocí transpozic sousedních prvků, což se opakovaně ukázalo jako velmi účinný způsob rozbíjení symetrií. Navíc je v porovnání s množinou \texttt{Transpositions} kompromisem mezi velikostí množiny a účinností rozbíjení symetrií. Množina \texttt{Transpositions} sice je schopna rozbít více symetrií, ale kvůli své velikosti je pro většinu zadání prakticky nepoužitelná.

Dalším opakovaným pozorováním bylo, že generující množiny permutací dosahovaly téměř ve všech experimentech lepších výsledků než stejně velké negenerující množiny. Tento trend se projevil při porovnání množin \texttt{Randtrans} a \texttt{Nongen}. Přestože experimentální výsledky samy o sobě nemohou tuto vlastnost dokázat, naznačují, že skutečnost, že množina permutací generuje celou grupu symetrií, může mít pozitivní vliv na účinnost rozbíjení symetrií a mělo by smysl se tímto v budoucnu teoreticky zabývat.

Významný vliv na výsledky měl také počet použitých permutací. Obecně platilo, že větší množiny permutací vedly k účinnějšímu rozbití symetrií, zároveň však zvyšovaly čas generování klauzulí i velikost výsledné SAT formule. Ukázalo se proto, že při návrhu množiny permutací je nutné hledat kompromis mezi její velikostí a dosaženou účinností.

Dalším důležitým parametrem je hodnota $k_{max}$, určující hloubku generování \texttt{Lex-Leader} podmínek. Experimenty ukázaly, že nejvyšší hodnota $k_{max}$ nemusí vést k nejkratším časům řešení. Přestože větší hodnoty odstraní více symetrií, současně vytvářejí více pomocných proměnných a klauzulí, jejichž zpracování může SAT řešič zpomalit. Optimální hodnota $k_{max}$ se navíc lišila podle typu řešeného problému.

Samostatnou pozornost si zaslouží porovnání s nástrojem \texttt{satsuma}. Ukázalo se, že množiny generátorů nalezené tímto nástrojem jsou i při relativně malém počtu permutací velmi kvalitní. V některých případech dosahoval nástroj \texttt{satsuma} lepších výsledků než náš přístup, v jiných naopak vedla specializovaná množina permutací navržená přímo pro daný problém k výrazně lepším výsledkům. Důležité je zmínit, že při použití druhé varianty nástroje \texttt{satsuma} byly výsledky velmi rozdílné. Varianta \texttt{fix} oproti variantě \texttt{lex} dosahovala u některých problémů dramaticky lepších výsledků. Pro stejnou množinu permutací generuje klauzule složitějším způsobem, díky kterému rozbíjí symetrie účinněji. Například u problému N královen nebo problému nalezení matice s řádky a sloupci se zadanými součty byly ale výsledky srovnatelné.

Celkově lze konstatovat, že navržený nástroj představuje univerzální přístup k rozbíjení symetrií SAT formulí založený na volitelných množinách permutací. Experimentální výsledky potvrzují, že vhodná volba množiny permutací i parametru $k_{max}$ může mít zásadní vliv na výslednou efektivitu. Současně experimenty odhalily několik zajímavých vlastností, zejména vliv generujících množin permutací a kvalitu množin nalezených nástrojem \texttt{satsuma}, které si zaslouží podrobnější teoretické zkoumání v budoucí práci.